\documentclass[11pt]{amsart}
\usepackage{amsmath,amssymb,amsthm}
\usepackage[margin=1.1in]{geometry}
\usepackage{hyperref}
\emergencystretch=2em

\newtheorem{theorem}{Theorem}
\newtheorem{lemma}[theorem]{Lemma}
\newtheorem{corollary}[theorem]{Corollary}
\newtheorem{proposition}[theorem]{Proposition}
\theoremstyle{remark}
\newtheorem{remark}[theorem]{Remark}

\newcommand{\Inf}{\mathrm{Inf}}
\newcommand{\Var}{\mathrm{Var}}
\newcommand{\E}{\mathbb{E}}

\title{The Fourier entropy of linear monotone DNFs on few variables}
\author{Yitzchak Shmalo}
\date{}

\begin{document}

\begin{abstract}
A monotone DNF is linear if any two of its terms share at most one variable. Over all linear monotone DNFs with at most
eight relevant variables, of which there are $552{,}251$ isomorphism classes on exactly eight, we determine by
exhaustive enumeration with certified logarithms the largest value of the ratio $H(f)/I(f)$ of Fourier entropy to total
influence and the largest value of $\min_iE_i(f)/\Inf_i(f)$, where $E_i$ is the erasure entropy at $i$. They are
$3.8891054198\ldots$ and $3.3047013224\ldots$, and each is attained by a single class. A transversal design on $21$
variables has $H/I=4.1346\ldots$, so neither value is the supremum over the class. We then give three sufficient
conditions for a coordinate with $E_i\le C\,\Inf_i$ in a linear monotone DNF: a load certificate that assumes no bound
on widths or degrees, a criterion that involves the loads of the variables of a single term only, and a bound in terms
of the bias of $f$. On the affine transversal designs with $2^k$ points per line and $k+2\lceil\log_2k\rceil+1$ lines
per point class, all of these conditions fail at every coordinate, as do the width and variance bounds of earlier work,
while $\Pr[F=0]$ tends to zero more slowly than any exponential in the number of variables.
\end{abstract}

\maketitle

\section{Introduction}

For $f\colon\{0,1\}^n\to\{-1,1\}$ write $f=\sum_S\hat f(S)\chi_S$ with $\chi_S(x)=(-1)^{\sum_{j\in S}x_j}$, and let $U$
be a random subset with $\Pr[U=S]=\hat f(S)^2$. The Fourier entropy is $H(f)=H(U)$, the influence of $i$ is
$\Inf_i(f)=\Pr[i\in U]$, and $I(f)=\sum_i\Inf_i(f)$. The Fourier entropy-influence (FEI) conjecture of Friedgut and
Kalai~\cite{FK96} asks for an absolute $C$ with $H(f)\le C\,I(f)$ for every Boolean $f$. It is known for symmetric
functions~\cite{OWZ}, for read-once formulas over gates of bounded arity~\cite{OT13} and over AND and OR
gates~\cite{CKLS}, and for read-$k$ decision trees with $k$ fixed~\cite{WWW}; its min-entropy version is known for
read-$k$ DNFs with $k$ fixed~\cite{Shalev,ACKSW}. The erasure entropy of $f$ at $i$ is
\[
E_i(f)=H\bigl(\mathbf 1[i\in U]\bigm| U\setminus\{i\}\bigr),
\]
and the FEI conjecture holds if and only if there is an absolute $c$ such that every non-constant Boolean function has
a coordinate with $E_i(f)\le c\,\Inf_i(f)$~\cite{Shm1}.

Let $F$ be a monotone DNF and $f=1-2F$. We call $F$ \emph{linear} if any two of its terms share at most one variable.
Read-once DNFs are the linear DNFs in which every variable lies in one term. For linear DNFs whose terms all have width
$w$ and in which every variable lies in at most $D\le2^{(1-\eta)(w-1)}$ terms, $E_i/\Inf_i=2+o(1)$ at every coordinate
as $w\to\infty$~\cite{Shm6}, and a DNF of width $w$ has a coordinate with $E_i/\Inf_i\le w+3+\log_2e$~\cite{Shm3};
for a monotone DNF one term of width $w$ is enough (Section~\ref{sec:prelim}). This paper concerns linear monotone DNFs with no restriction on widths or degrees. Since the minimal DNF
of a monotone function is unique, linearity is a property of the function.

\subsection{Small dimensions}
Call two functions equivalent if they differ by a permutation of the variables; $H$, $I$ and the multiset of pairs
$(E_i,\Inf_i)$ are invariants.

\begin{theorem}\label{thm:census}
For $n\le8$, Table~\ref{tab:census} lists the number of equivalence classes of linear monotone functions with exactly $n$
relevant variables, the number of such functions on a fixed set of $n$ variables, and the largest values of $H/I$ and of
$\min_iE_i/\Inf_i$ over them. Over all linear monotone functions with at most eight relevant variables,
\[
3.30470132244542\le\max\min_i\frac{E_i(f)}{\Inf_i(f)}\le3.30470132245649,
\]
attained only by $F=y_1y_2\vee y_2y_3\vee y_3y_4\vee y_4y_1\vee z_1z_2\vee z_2z_3\vee z_3z_4\vee z_4z_1$, and
\[
3.88910541981418\le\max\frac{H(f)}{I(f)}\le3.88910541981709,
\]
attained only by $F=x_1\vee x_2\vee x_3\vee x_4\vee y_1y_2\vee y_2y_3\vee y_3y_4\vee y_4y_1$.
\end{theorem}

\begin{table}[h]
\centering
\begin{tabular}{rrrcc}
$n$ & classes & functions & $\max\min_iE_i/\Inf_i$ & $\max H/I$\\\hline
2 & 2 & 2 & 2.0000000000 & 2.0000000000\\
3 & 5 & 9 & 2.6724889840 & 2.9558895823\\
4 & 16 & 97 & 2.6340033909 & 3.4024755512\\
5 & 72 & 2{,}625 & 2.6958875202 & 3.6237403997\\
6 & 592 & 185{,}521 & 2.9444953039 & 3.7397640179\\
7 & 10{,}808 & 35{,}406{,}319 & 2.9546838142 & 3.8043570698\\
8 & 552{,}251 & 18{,}785{,}639{,}553 & 3.3047013225 & 3.8891054198\\
\end{tabular}
\caption{Linear monotone functions with exactly $n$ relevant variables. Values rounded to ten decimals; at $n=8$ the
second largest values are $3.1084496901$ and $3.8663646986$.}\label{tab:census}
\end{table}

As far as we know this is the first exhaustive computation of the FEI ratio on a class of this size. Biswas and
Sarkar~\cite{BS} searched all Boolean functions on at most five variables for the min-entropy version of the
conjecture; the largest ratio of min-entropy to influence they found there is $16/7$.

The maxima of Theorem~\ref{thm:census} are not suprema over the class. The function $\mathrm{AND}_n$ is a linear DNF
with one term, and $H/I\to4$ along it. The transversal designs of Section~\ref{sec:td}, which have many intersecting
terms of equal width, give larger values: on $21$ variables, with $49$ terms of width $3$ and every variable in $7$ of
them, $H/I=4.1345862211$ (Table~\ref{tab:td}). On these designs $\min_iE_i/\Inf_i$ stays between $2.10$ and $2.72$.
All these values lie well below the best known lower bound $C>6.454784$ on the constant of the conjecture, due to
Hod~\cite{Hod}, so they say nothing about $C$ itself.

\subsection{Certificates and an obstruction}
For a variable $j$ put $\lambda_j=\sum_{t\ni j}2^{-(|t|-1)}$, the sum over the terms through $j$; we call it the
\emph{load} of $j$. In a term of width $L+1$ that meets no other term every variable has load $2^{-L}$.

\begin{theorem}\label{thm:local}
Let $F$ be a linear monotone DNF and $t$ a term of width $L+1$ with $L\ge3$. If every variable $j$ of $t$ satisfies
\[
\lambda_j\le\Bigl(1+\frac1{64L^2}\Bigr)2^{-L},
\]
then $E_i(f)<6\,\Inf_i(f)$ for every $i\in t$.
\end{theorem}

Nothing is assumed about the variables outside $t$, and $t$ need not be a shortest term. Theorem~\ref{thm:local} is a
proved by a variant of the argument for a finite certificate, Proposition~\ref{prop:mix}, which bounds $|E_i/\Inf_i-2|$ by a quantity computed
from the terms through $i$ and the loads of their variables, with no maximum width or maximum degree. The third
condition is a bound by the bias of $f$ (Lemma~\ref{lem:bias}): if $m=\min(\Pr[f=1],\Pr[f=-1])>0$ then
$H(f)/I(f)<2+2n/\log_2(1/m)$, for every Boolean $f$ on $n$ variables.

Section~\ref{sec:td} shows that these conditions do not combine into a proof of the FEI conjecture for linear monotone
DNFs. For $Q$ a prime power and $w\le Q$ let $\mathrm{TD}(w,Q)$ be the DNF whose variables are the points of
$A\times\mathbb F_Q$ for a $w$-set $A\subseteq\mathbb F_Q$ and whose terms are the graphs of the affine maps
$x\mapsto ax+b$ restricted to $A$.

\begin{theorem}\label{thm:td}
Let $Q=2^k$ and $w=k+2\lceil\log_2k\rceil+1$, and let $F_k=\mathrm{TD}(w,Q)$. As $k\to\infty$:
\begin{enumerate}
\item every term has width $w\to\infty$, and the hypothesis of Theorem~\ref{thm:local} fails for every term;
\item the right side of the certificate of Proposition~\ref{prop:mix} tends to infinity at every coordinate,
uniformly over all choices of its parameters;
\item $\Pr[F_k=0]\to0$ but $\log_2(1/\Pr[F_k=0])/n\to0$, so the bias bound tends to infinity;
\item the variance bound of Lemma~\ref{lem:rel} tends to infinity at every coordinate.
\end{enumerate}
\end{theorem}

Theorem~\ref{thm:td} says nothing about the actual ratios on $F_k$, which we can compute only for small designs; there
they are small (Table~\ref{tab:td}).

\section{Preliminaries}\label{sec:prelim}

We use the notation of~\cite{Shm6}. Fix $i$ and write $f^{(b)}$ for the restriction $x_i=b$. Put
$g=\tfrac12(f^{(0)}+f^{(1)})$ and $h=\tfrac12(f^{(0)}-f^{(1)})$, so that $\hat g(T)=\hat f(T)$ and $\hat
h(T)=\hat f(T\cup\{i\})$ for $T\subseteq[n]\setminus\{i\}$, and $\E h^2=\Inf_i(f)=:\delta$. With $p(T)=\hat g(T)^2$,
$q(T)=\hat h(T)^2$ and
\[
\Phi(a,b)=\eta(a)+\eta(b)-\eta(a+b),\qquad\eta(x)=-x\log_2x,
\]
we have $E_i(f)=\sum_T\Phi(p(T),q(T))$. The function $\Phi$ is concave and homogeneous of degree one, hence
superadditive, it is increasing in each argument, and
\begin{equation}\label{eq:phiq}
\Phi(a,b)\le b\log_2\Bigl(e\,\frac{a+b}{b}\Bigr).
\end{equation}
If $F$ is monotone and $f=1-2F$, let $N$ and $M$ be the indicators that $F$ vanishes at $x_i=0$ and at $x_i=1$. Then
$h=N-M$ is the indicator that $i$ is pivotal, $\hat h(\emptyset)=\delta$, and $\hat g(T)=2\hat N(T)-\hat h(T)$ for
$T\ne\emptyset$. We use the separation inequality of~\cite{Shm6}.

\begin{proposition}[{\cite[Proposition~2]{Shm6}}]\label{prop:sep}
Let $F$ be monotone, $f=1-2F$, $\delta=\Inf_i(f)>0$. Let $\Omega$ be a family of non-empty subsets of
$[n]\setminus\{i\}$ and let $\varepsilon\in(0,1]$ and $\kappa\ge0$ satisfy $\sum_{T\notin\Omega}\hat h(T)^2\le
\varepsilon\delta$ and $\sum_{T\in\Omega}\hat N(T)^2\le\kappa\delta$. Then
\[
2-2\varepsilon-\frac{8}{\ln2}\sqrt\kappa\;\le\;\frac{E_i(f)}{\delta}\;\le\;2+\frac{8}{\ln2}\sqrt\kappa
+\varepsilon\log_2\Bigl(e\Bigl[3+\frac{8\Var N}{\varepsilon\delta}\Bigr]\Bigr)+\delta\log_2\bigl(e[1+\delta^{-2}]\bigr).
\]
\end{proposition}

The next bound improves the variance lemma of~\cite{Shm3} for monotone functions.

\begin{lemma}\label{lem:rel}
Let $F$ be monotone, $f=1-2F$, $v=\Var f$ and $\delta=\Inf_i(f)>0$. Then
\[
\frac{E_i(f)}{\delta}\le2|\E f|+\log_2\frac{e\,v}{\delta}.
\]
\end{lemma}

\begin{proof}
First, $\Phi(a,b)\le2\sqrt{ab}$. By symmetry and homogeneity it suffices to take $b=1$ and $a=x^2$ with $x\ge1$. Put
$\ell(x)=2x\ln2-(x^2+1)\ln(x^2+1)+x^2\ln(x^2)$, which is $\ln2$ times $2x-\Phi(x^2,1)$. Then $\ell(1)=0$ and
$\ell'(x)=2\ln2-2x\ln(1+x^{-2})$. The function $x\ln(1+x^{-2})$ decreases on $[1,\infty)$: its derivative is
$\ln(1+y)-2y/(1+y)$ with $y=x^{-2}\in(0,1]$, which is non-positive because it vanishes at $y=0$ and has derivative
$-(1-y)/(1+y)^2\le0$. So $\ell'\ge\ell'(1)=0$ and $\ell\ge0$.

At $T=\emptyset$, $p=(\E f)^2$ and $q=\delta^2$, so the term is at most $2|\E f|\delta$. By Parseval the non-empty
sets carry $P=\sum_{T\neq\emptyset}p(T)=v-\delta$ and $Q=\sum_{T\ne\emptyset}q(T)=\delta-\delta^2$. By superadditivity
and~\eqref{eq:phiq},
\[
\sum_{T\ne\emptyset}\Phi(p,q)\le\Phi(P,Q)\le Q\log_2\frac{e(v-\delta^2)}{Q}\le Q\log_2\frac{ev}Q\le\delta\log_2\frac{ev}\delta,
\]
since $x\mapsto x\log_2(ev/x)$ increases on $(0,v]$ and $Q\le\delta\le v$; here $\delta\le v$ because
$\delta\le2\mu\le4\mu(1-\mu)=v$ with $\mu=\min(\Pr[f=1],\Pr[f=-1])\le1/2$. If $\delta=1$ the non-empty part vanishes.
\end{proof}

For comparison, \cite{Shm3} has $2+\log_2e+\log_2(1+v/\delta)$ for every Boolean $f$. The proof of the width bound
of~\cite{Shm3} uses a single term: if $t$ is any term of width $w$ of a monotone DNF, then some $i\in t$ has
$\Inf_i(f)\ge2\Pr[F=0]/(2^w-1)$, and Lemma~\ref{lem:rel} gives $E_i/\Inf_i\le2+\log_2(2e(2^w-1))$ there.

\section{The census}\label{sec:census}

\subsection{Enumeration}
A linear monotone function is determined by its minimal DNF, a linear antichain of subsets of the variables, and the
function has exactly $n$ relevant variables when the union of the terms is $[n]$. Singleton terms share no variable
with any other term. It therefore suffices to enumerate, up to isomorphism, the linear hypergraphs on $[n]$ with edges of
size at least two and isolated vertices allowed, and to add a singleton edge at every isolated vertex.

The enumeration processes widths from $n$ down to $2$. At width $w$ it repeatedly adds to each family every $w$-set
that meets each existing edge in at most one vertex, keeping one family per isomorphism class, until no new class
appears, and then passes to $w-1$. Every linear hypergraph is reached: list its edges by decreasing width, and its
prefixes are linear hypergraphs reached at the earlier stages. Isomorphism is decided by an exact canonical form. The
vertices are partitioned by the multiset of the sizes of their incident edges, the partition is refined, and
unresolved cells are individualised; at each leaf the relabelled edge list is sorted, and the least leaf is the
canonical form. Two families with the same canonical form are related by the permutation read off the two leaves.

\subsection{Completeness}
Two independent counts confirm that no class is missing or repeated. The number of labelled linear hypergraphs on $[n]$
with all edges of size at least two, isolated vertices allowed, is computed by a recursion on pairs: choose an uncovered pair $uv$; either no edge
of size at least three contains it, and then $uv$ is either an edge or not, or exactly one such edge does, and its
pairs are removed. This is compared with $\sum n!/|\mathrm{Aut}(H)|$ over the census, with $|\mathrm{Aut}(H)|$ counted
at the canonical leaves. Adding a singleton term at each isolated vertex is a bijection onto the functions with
exactly $n$ relevant variables, so both counts are the third column of Table~\ref{tab:census}, and they agree for every
$n\le8$. The census for $n\le5$ was also compared with a brute-force enumeration of all monotone functions.

A second program, written separately and sharing no code with the first, reads the census and checks for each family
that the terms are distinct, pairwise meet in at most one variable, form an antichain and cover $[n]$; it computes a
canonical form by a different method (the least sorted relabelled edge list over all relabellings that preserve the
vertex invariant) and the automorphism count, and recomputes both ratios. It finds no invalid family and no repeated
class, reproduces the third column of Table~\ref{tab:census}, and agrees with the certified values below to the
precision of double arithmetic.

\subsection{Certified values}
For each function the Walsh transform of $f=1-2F$ is computed in integers, so every $p(T)$ and $q(T)$ is an exact
dyadic rational with denominator $4^n$, and $H$, $I$ and every $E_i$ are integer combinations of $\log_2k$ for
integers $k\le4^n$, divided by $4^n$. Enclosures of $\log_2k$ are computed in integer arithmetic from
$\ln x=2\sum_{j\ge0}z^{2j+1}/(2j+1)$ with $z=(x-1)/(x+1)\le1/3$ and a geometric bound on the tail, rounded outwards to
units of $2^{-52}$; sums are accumulated in signed $128$-bit integers and every floating conversion and division is
followed by an outward step. This encloses $H/I$ and every $E_i/\Inf_i$ for every function in the census, and the
intervals in Theorem~\ref{thm:census} are the resulting enclosures for the two maximisers. In both cases the upper end
of every other class's enclosure lies below the lower end of the maximiser's, which proves uniqueness. The exact values
are also available as rational combinations of logarithms of primes.

\subsection{The maximisers}
The maximiser of $\min_iE_i/\Inf_i$ is the OR of two copies, on disjoint variables, of the cycle DNF
$C_4=y_1y_2\vee y_2y_3\vee y_3y_4\vee y_4y_1$, for which $H/I=3.6125175459$. The maximiser of $H/I$ is the OR of
$\mathrm{OR}_4$ and $C_4$ on disjoint variables. Both maxima grow at every step from $n=4$ to $n=8$, so the census
gives no indication of the suprema.

\section{The load certificate}\label{sec:mix}

Let $F$ be a linear monotone DNF and $i$ a variable that lies in no singleton term. Let $t$ range over the terms
through $i$, put $B_t=t\setminus\{i\}$, $L_t=|B_t|\ge1$, $a_t=2^{-L_t}$, and $b=\sum_ta_t=\lambda_i$. By linearity the
sets $B_t$ are pairwise disjoint. A \emph{background term} is a term not containing $i$. For a variable $j\ne i$ put
\[
\ell_j=\sum_{s\ni j,\ i\notin s}2^{-(|s|-1)},\qquad \eta_t=\sum_{j\in B_t}\ell_j,\qquad
\theta=1-\frac1b\Bigl[\sum_ta_t\eta_t+\sum_{s<t}a_sa_t\Bigr].
\]
For $j\in B_t$, a term through $j$ other than $t$ does not contain $i$, since it would share two variables with $t$;
so $\ell_j=\lambda_j-a_t$. For integers $1\le k_t\le L_t$ put
\[
\gamma_t=a_t\sum_{r<k_t}\binom{L_t}r,\qquad
\alpha_t=2^{\binom{k_t}2-2k_t}\sum_{J\subseteq B_t,\ |J|=k_t}\ \frac{\prod_{j\in J}\ell_j}{1-\frac12\sum_{j\in J}\ell_j},
\]
and
\[
\varepsilon=\min\Bigl(1,\sum_t\gamma_t\Bigr),\qquad \kappa=\frac{\sum_t\alpha_t}{b\theta}.
\]

\begin{proposition}\label{prop:mix}
If $\theta>0$, every denominator in $\alpha_t$ is positive, and $b\le1/4$, then
\[
\Bigl|\frac{E_i(f)}{\Inf_i(f)}-2\Bigr|\le2\varepsilon+\frac8{\ln2}\sqrt\kappa
+\varepsilon\log_2\Bigl(e\Bigl[3+\frac8{\varepsilon b\theta}\Bigr]\Bigr)+b\log_2\bigl(e[1+b^{-2}]\bigr).
\]
\end{proposition}

When all terms have width $w$ and all loads are bounded through the degree, this is~\cite[Theorem~6]{Shm6}; the proof
follows the three lemmas there, and we indicate the changes. Let $N$ be the indicator that every background term fails,
$A=\E N$, $A_t$ the indicator that all of $B_t$ equal $1$, and $h=N\cdot\mathbf 1[\bigvee_tA_t]$, so that
$\delta=\E h=\Inf_i(f)$.

\begin{lemma}\label{lem:m1}
$Ab\theta\le\delta\le Ab$.
\end{lemma}

\begin{proof}
The upper bound is Harris's inequality for the decreasing $N$ and the increasing $A_t$, and the union bound. For the
lower bound fix $t$, let $R_t$ indicate that every background term disjoint from $B_t$ fails, and let $V_t$ be the OR
of the residual conjunctions $s\setminus B_t$ of the background terms $s$ meeting $B_t$. Each such $s$ meets $B_t$ in one
variable, so $\Pr[V_t]\le\eta_t$, and $R_t$ and $V_t$ do not depend on $B_t$. By Harris,
$\Pr[N,A_t]=a_t\Pr[R_t,\overline{V_t}]\ge a_t\Pr[R_t](1-\eta_t)\ge a_tA(1-\eta_t)$; if $\eta_t>1$ this holds
trivially. For $s\ne t$, $\Pr[N,A_s,A_t]\le Aa_sa_t$, since $B_s$ and $B_t$ are disjoint. Bonferroni's inequality gives
$\delta\ge A\sum_ta_t(1-\eta_t)-A\sum_{s<t}a_sa_t=Ab\theta$.
\end{proof}

Let $\Omega$ be the family of sets $T\subseteq[n]\setminus\{i\}$ with $|T\cap B_t|\ge k_t$ for some $t$.

\begin{lemma}\label{lem:m2}
$\sum_{T\notin\Omega}\hat h(T)^2\le\varepsilon\delta$.
\end{lemma}

\begin{proof}
Order the terms through $i$ and put $h_t=NA_t\prod_{u<t}(1-A_u)$. These have disjoint supports and sum to $h$, so
$\sum_t\|h_t\|_2^2=\delta$. As in~\cite[Lemma~4]{Shm6}, $h_t=A_t\Psi_t$ with $\Psi_t$ independent of $B_t$, so the
weight of $\hat h_t^2$ on sets meeting $B_t$ in exactly $r$ variables is $\binom{L_t}ra_t\|h_t\|_2^2$, and the weight
of $\hat h_t^2$ off $\Omega$ is at most $\gamma_t\|h_t\|_2^2$. By the triangle inequality and Cauchy--Schwarz,
\[
\Bigl(\sum_{T\notin\Omega}\hat h(T)^2\Bigr)^{1/2}\le\sum_t\gamma_t^{1/2}\|h_t\|_2\le\Bigl(\sum_t\gamma_t\Bigr)^{1/2}\delta^{1/2}.
\]
The bound by $\delta$ is trivial.
\end{proof}

\begin{lemma}\label{lem:m3}
$\sum_{T\in\Omega}\hat N(T)^2\le A\sum_t\alpha_t$.
\end{lemma}

\begin{proof}
Every $T\in\Omega$ contains a $k_t$-subset $J$ of some $B_t$, and $\sum_{T\supseteq J}\hat N(T)^2=\|D_JN\|_2^2$. Fix
such $J$, $k=k_t$. A background term meets $J$ in at most one variable. Let $R_J$ indicate that every background term
avoiding $J$ fails, and $Q_j$, for $j\in J$, the OR of the residual conjunctions $s\setminus\{j\}$ of the background terms
$s\ni j$; then $\Pr[Q_j]\le\ell_j$, and as in~\cite[Lemma~5]{Shm6}, $\|D_JN\|_2^2=4^{-k}\Pr[R_J,\bigwedge_jQ_j]\le
4^{-k}\Pr[R_J]\Pr[\bigwedge_jQ_j]$. For $\bigwedge_jQ_j$ a background term $s_j\ni j$ must have its residual conjunction
satisfied for each $j$. The $s_j$ are distinct, two of the residual sets share at most one variable, so their union has
at least $\sum_j(|s_j|-1)-\binom k2$ variables, and summing over the choices gives $\Pr[\bigwedge_jQ_j]\le2^{\binom
k2}\prod_{j\in J}\ell_j$. Finally $A\ge\Pr[R_J]-\sum_j\frac12\Pr[R_J,Q_j]\ge\Pr[R_J](1-\frac12\sum_j\ell_j)$ by Harris,
since $x_j$ is independent of $R_J$ and $Q_j$. Hence $\|D_JN\|_2^2$ is at most $A$ times the summand of $\alpha_t$
indexed by $J$; sum over $J$ and $t$.
\end{proof}

\begin{proof}[Proof of Proposition~\ref{prop:mix}]
By Lemmas~\ref{lem:m3} and~\ref{lem:m1}, $\sum_\Omega\hat N^2\le A\sum_t\alpha_t\le\kappa\delta$, and by
Lemma~\ref{lem:m2} the first hypothesis of Proposition~\ref{prop:sep} holds with $\varepsilon$; note $\varepsilon\ge
b>0$ because $\gamma_t\ge a_t$. The sets in $\Omega$ are non-empty. In the upper bound of Proposition~\ref{prop:sep},
$\Var N\le A\le\delta/(b\theta)$, and $\delta\le Ab\le b\le1/4$; the function $x\log_2(e[1+x^{-2}])$ is increasing on
$(0,1/4]$~\cite{Shm6}. The lower bound of Proposition~\ref{prop:sep} is within the stated error.
\end{proof}

\section{A criterion on one term}\label{sec:local}

\begin{proof}[Proof of Theorem~\ref{thm:local}]
Put $a=2^{-L}$, $u=1/(64L^2)$, $\theta_0=1-Lua$, and fix $i\in t$. The terms through $i$ other than $t$ have total
weight $\lambda_i-a\le ua$, and every $j\in t\setminus\{i\}$ has $\ell_j=\lambda_j-a\le ua$. Let $N$, $A$, $A_t$ and
$h$ be as in Section~\ref{sec:mix}, and split $h=h_t+h'$ with $h_t=NA_t$, $h'=N(1-A_t)\mathbf 1[\bigvee_{s\ne t}A_s]$,
which have disjoint supports. By Harris and the proof of Lemma~\ref{lem:m1},
\[
\|h_t\|_2^2\le Aa,\qquad\|h'\|_2^2\le Aua,\qquad Aa\theta_0\le\delta\le A(1+u)a.
\]
Let $\Omega$ be the family of sets that contain at least one variable of $t\setminus\{i\}$. The weight of $\hat h_t^2$
off $\Omega$ is $a\|h_t\|_2^2\le Aa^2$, so
\[
\Bigl(\sum_{T\notin\Omega}\hat h(T)^2\Bigr)^{1/2}\le\sqrt{Aa}\,(\sqrt a+\sqrt u),\qquad
\sum_{T\notin\Omega}\hat h(T)^2\le\varepsilon_0\delta,\quad\varepsilon_0=\frac{(\sqrt a+\sqrt u)^2}{\theta_0}.
\]
The proof of Lemma~\ref{lem:m3} with $k=1$ for the single block $t\setminus\{i\}$ gives $\sum_\Omega\hat N^2\le
A\,Lua/(4(1-ua/2))\le\kappa_0\delta$ with $\kappa_0=Lu/(4\theta_0(1-ua/2))$. Proposition~\ref{prop:sep}, with $\Var
N\le A\le\delta/(a\theta_0)$ and $\delta\le b_0=(1+u)a$, gives
\[
\frac{E_i}{\delta}\le2+\frac8{\ln2}\sqrt{\kappa_0}+\varepsilon_0\log_2\Bigl(e\Bigl[3+\frac{8}{\varepsilon_0a\theta_0}\Bigr]\Bigr)
+b_0\log_2\bigl(e[1+b_0^{-2}]\bigr).
\]
Now let $L\ge3$, so $a\le1/8$ and $\sqrt u=1/(8L)\le1/24$. Then $\theta_0\ge1-1/1536$, $\varepsilon_0<0.157$, and
$L\varepsilon_0<0.47$, since $La$, $L\sqrt{au}=2^{-L/2}/8$ and $Lu$ decrease in $L$. Next
$\kappa_0<1/(256L\cdot0.999)<0.0014$, so the second term is below $11.55\cdot0.0374<0.44$. Since
$\varepsilon_0a\theta_0\le1$, the third term is at most
\[
\varepsilon_0\Bigl[\log_2\frac{11e}{\theta_0}+L+\log_2\frac1{\varepsilon_0}\Bigr]<0.157\cdot5+0.47+\frac1{e\ln2}<1.79.
\]
Finally $b_0<0.126$, and since $x\log_2(e[1+x^{-2}])$ increases on $(0,1/4]$ the last term is below
$0.126\log_2(e\cdot64)<0.95$. Altogether $E_i/\delta<2+0.44+1.79+0.95<6$.
\end{proof}

For terms of width at most three the bound $2+\log_2(2e(2^w-1))<7.3$ of Section~\ref{sec:prelim} holds at a variable
of largest influence in the term, with no hypothesis on loads.

\section{Bias}\label{sec:bias}

\begin{lemma}\label{lem:bias}
Let $f\colon\{0,1\}^n\to\{-1,1\}$, $m=\min(\Pr[f=1],\Pr[f=-1])>0$ and $r=\log_2(1/m)$. Then
\[
\frac{H(f)}{I(f)}\le2+\frac{2(n-2+\log_2e)}r.
\]
\end{lemma}

\begin{proof}
Let $v=1-\hat f(\emptyset)^2=4m(1-m)$ be the weight off the empty set. Conditioning on whether $U$ is empty,
$H(f)\le h_2(v)+v\log_2(2^n-1)\le v\log_2(e/v)+vn$, using $-(1-v)\log_2(1-v)\le v\log_2e$. The right side increases
in $v$ on $(0,2^n]$ (its derivative is $n-\log_2v$), and $v\le4m\le2$, so $H(f)\le4m(n+r+\log_2e-2)$. A set $S$ of $s$ vertices of the $n$-cube spans at most
$\frac12s\log_2s$ edges~\cite{Harper,Hart}, so at least $s(n-\log_2s)$ edges leave it; with $S$ the minority set,
$s=m2^n$, this gives $I(f)\ge2\cdot2^{-n}s(n-\log_2s)=2mr$. Divide.
\end{proof}

For $\mathrm{AND}_n$ one has $m=2^{-n}$, $r=n$, and both sides tend to $4$. The bound is useful only when $m$ is
exponentially small in $n$: if $m\le2^{-\beta n}$ it gives $H/I<2+2/\beta$.

\section{Transversal designs}\label{sec:td}

Let $Q$ be a prime power, $A\subseteq\mathbb F_Q$ with $|A|=w$, and let $\mathrm{TD}(w,Q)$ have variables $(x,y)\in
A\times\mathbb F_Q$ and terms $t_{a,b}=\{(x,ax+b):x\in A\}$ for $a,b\in\mathbb F_Q$. It has $n=wQ$ variables and $Q^2$
terms of width $w$; the point $(x,y)$ lies on the $Q$ terms with $b=y-ax$; two distinct affine maps agree at most once,
so the DNF is linear. Every variable has load $\lambda=Q2^{-(w-1)}$, and every variable of a term $t$ through $i$, other
than $i$, has $\ell_j=(Q-1)2^{-(w-1)}$.

\begin{table}[h]
\centering
\begin{tabular}{rrrrcccc}
$w$ & $Q$ & $n$ & terms & $\Pr[F=0]$ & $H/I$ & $\min_iE_i/\Inf_i$ & $\max_iE_i/\Inf_i$\\\hline
3 & 3 & 9 & 9 & 0.4180 & 3.0856094956 & 2.1015004988 & 2.1015004988\\
3 & 5 & 15 & 25 & 0.1435 & 3.8797848238 & 2.7150952652 & 2.7150952652\\
4 & 5 & 20 & 25 & 0.3219 & 3.9024430124 & 2.4713311010 & 2.4713311010\\
3 & 7 & 21 & 49 & 0.0408 & 4.1345862211 & 2.5335124504 & 2.5335124504\\
5 & 5 & 25 & 25 & 0.5332 & 3.8799716886 & 2.2003373077 & 2.2003373077\\
\end{tabular}
\caption{$\mathrm{TD}(w,Q)$ with $A=\{0,\dots,w-1\}$, from exact Walsh transforms.}\label{tab:td}
\end{table}

\begin{proof}[Proof of Theorem~\ref{thm:td}]
Write $L=w-1=k+2\lceil\log_2k\rceil$, $a=2^{-L}$, $\lambda=Qa=2^{-2\lceil\log_2k\rceil}\in[1/(4k^2),1/k^2]$ and
$\ell=(Q-1)a<\lambda$; note $w\le Q$ for $k\ge3$ and $L/k\to1$.

(1) $\lambda_j/a=Q$, far above $1+1/(64L^2)$.

(2) Fix $i$. It lies in $d=Q$ terms with $L_t=L$, so $b=\lambda$, $\eta_t=L\ell$ and
$\theta=1-L\ell-(Q-1)a/2=1-O(1/k)$. Choose thresholds $k_t$. If at least $Q/2$ of them are at least $4$, then
$\sum_t\gamma_t\ge\frac Q2a\binom L3=\frac\lambda2\binom L3\ge\frac{L^3(1-o(1))}{48k^2}\to\infty$, so eventually
$\varepsilon=1$ and the third term of Proposition~\ref{prop:mix} is at least $\log_2(8e/\lambda)\to\infty$.
Otherwise at least $Q/2$ of them lie in $\{1,2,3\}$, and a block with $k_t=r$ has
$\alpha_t=\binom Lr2^{\binom r2-2r}\ell^r/(1-r\ell/2)$, so
\[
\kappa\ge\frac{Q}{2}\cdot\frac{\min_{r\le3}\alpha^{(r)}}{Qa\theta}=\frac{\min_{r\le3}\alpha^{(r)}}{2a\theta}.
\]
Since $\ell\ge\lambda/2$ and $1/a=Q/\lambda$, $\alpha^{(r)}/a$ is of order $\binom Lr\lambda^{r-1}Q$, that is $LQ$,
$Q$ and $Q/L$ for $r=1,2,3$ up to bounded factors. So $\kappa\to\infty$.

(3) Let $Z$ be the number of satisfied terms and $p=2^{-w}$. Then $\E Z=Q^2p\ge2^k/(8k^2)\to\infty$. Terms with
disjoint variables are independent; two intersecting terms have covariance $2^{-(2w-1)}-p^2=p^2$, and each term meets
$w(Q-1)$ others. Hence $\Var Z/(\E Z)^2\le1/\E Z+w/Q\to0$, and $\Pr[F=0]=\Pr[Z=0]\to0$ by Chebyshev's inequality. On
the other hand the failures of the terms are decreasing events, so by Harris $\Pr[F=0]\ge(1-p)^{Q^2}$ and
\[
\frac{\log_2(1/\Pr[F=0])}{n}\le\frac Qw\log_2\frac1{1-p}=O\Bigl(\frac{2^{k-w}}w\Bigr)=O(k^{-3}).
\]
Once $\Pr[F=0]<1/2$ the minority value is $F=0$, so $m=\Pr[F=0]$, and the bound of Lemma~\ref{lem:bias}, which
exceeds $2(n-1)/\log_2(1/\Pr[F=0])$, tends to infinity.

(4) Fix $i$ and let $N$ and $A$ be as in Section~\ref{sec:mix}. Since $x_i$ is independent of $N$ and of the $A_t$,
$\Pr[F=0]=A-\frac12\delta$, and $\delta\le A\lambda$ by Lemma~\ref{lem:m1}. Hence $\delta/\Pr[F=0]\le\lambda/(1-\lambda/2)$
and $\Var f/\delta=4\Pr[F=0]\Pr[F=1]/\delta\ge4\Pr[F=1](1-\lambda/2)/\lambda\to\infty$, uniformly in $i$, so the right
side of Lemma~\ref{lem:rel} tends to infinity.
\end{proof}

For $\mathrm{TD}(w,2^k)$ the width bound of Section~\ref{sec:prelim} is of order $w$, and Theorem~\ref{thm:local} does
not apply to any term. So a linear monotone DNF in which every restriction has a coordinate covered by one of these
bounds satisfies $H\le C\,I$ by the induction of~\cite{Shm1}, but the designs show that not every linear monotone DNF
has such a coordinate even at the root. Covering them would need either a separation estimate that uses the loads
conditional on the background event $N$, rather than the unconditional loads, or a bound for averages over whole terms
in the spirit of decision-tree arguments. Table~\ref{tab:td} suggests that the actual ratios on the designs remain
small.

\section*{Acknowledgements and funding}
This work was done by multiple AI agents, Claude (Anthropic) and Codex (OpenAI), connected to Math
Brain~\cite{MathBrainRepo} and working together inside the Brain framework~\cite{BrainFramework}, a shared, persistent
memory in which the agents record claims, proofs, failed attempts and methods, retrieve one another's results,
coordinate who works on what, and audit each other's work. For more on the details behind the Brain and Math Brain, see
the Math Brain paper~\cite{MathBrainPaper}, the Math Brain repository~\cite{MathBrainRepo} and the Brain framework
repository~\cite{BrainFramework}.

The research presented in this paper was supported by the European Research Council (ERC) under the European Union's
Horizon Europe research and innovation programme (grant agreement No.~101041711), by the Simons Foundation as part of
the Collaboration on the Mathematical and Scientific Foundations of Deep Learning, by Heights Labs, by Convex Nexus
Capital, by the Israel Science Foundation (grant number 2258/19), and by the Israel Science Foundation (ISF Grant
4101/25). Computational resources were provided in part by Google Cloud through its education credits program (credit
530550103).

\section*{Declaration of competing interest}
The author holds equity in Heights Labs and in Convex Nexus Capital and has a management role at Convex Nexus Capital.
These entities are included in the funding declaration above.

\section*{Declaration of generative AI and AI-assisted technologies}
The paper comes out of Math Brain, an ongoing project on the use of AI systems for research mathematics, and its
results were developed with the Brain~\cite{BrainFramework} and Math Brain~\cite{MathBrainRepo}. In November 2026 a
Claude Fable agent (Anthropic) and Codex agents running OpenAI's GPT-6.1, connected to the Brain and working on the same
problem, studied erasure entropies of DNFs; a Codex agent wrote the census programs, ran the census and found the
certificates, the transversal-design obstruction and their proofs. Claude Opus 5.5 (Anthropic) rederived each proof,
wrote the independent census check and the design computations, prepared the manuscript and audited it. These checks
are not peer review. The author directed the work and is responsible for the content of the paper and for decisions
concerning its dissemination.

\section*{Code and data availability}
Code and outputs: \url{https://github.com/yspennstate/fei-linear-dnf-census}.
The paper is part of the Math Brain project, whose public repository is
\url{https://github.com/yspennstate/math-brain-public}~\cite{MathBrainRepo}.

\begin{thebibliography}{99}
\bibitem{ACKSW} S.~Arunachalam, S.~Chakraborty, M.~Kouck\'y, N.~Saurabh, R.~de~Wolf, Improved bounds on Fourier
entropy and min-entropy, ACM Trans. Comput. Theory 13 (2021), arXiv:1809.09819.
\bibitem{BS} A.~Biswas, P.~Sarkar, A lower bound on the constant in the Fourier min-entropy/influence conjecture,
Discrete Appl. Math. (2025), doi:10.1016/j.dam.2024.12.011, arXiv:2212.07713.
\bibitem{CKLS} S.~Chakraborty, R.~Kulkarni, S.~V.~Lokam, N.~Saurabh, Upper bounds on Fourier entropy, Theoret.
Comput. Sci. 654 (2016), 92--112.
\bibitem{FK96} E.~Friedgut, G.~Kalai, Every monotone graph property has a sharp threshold, Proc. Amer. Math. Soc.
124 (1996), 2993--3002.
\bibitem{Harper} L.~H.~Harper, Optimal assignments of numbers to vertices, J. Soc. Indust. Appl. Math. 12 (1964),
131--135.
\bibitem{Hart} S.~Hart, A note on the edges of the $n$-cube, Discrete Math. 14 (1976), 157--163.
\bibitem{Hod} R.~Hod, Improved lower bounds for the Fourier entropy/influence conjecture via lexicographic functions,
arXiv:1711.00762, 2017.
\bibitem{OT13} R.~O'Donnell, L.-Y.~Tan, A composition theorem for the Fourier entropy-influence conjecture, ICALP 2013.
\bibitem{OWZ} R.~O'Donnell, J.~Wright, Y.~Zhou, The Fourier entropy-influence conjecture for certain classes of
Boolean functions, ICALP 2011, Lecture Notes in Comput. Sci. 6755, 330--341.
\bibitem{Shalev} G.~Shalev, On the Fourier entropy influence conjecture for extremal classes, arXiv:1806.03646, 2018.
\bibitem{Shm1} Y.~Shmalo, The one-coordinate step in the Fourier entropy-influence conjecture,
\url{https://github.com/yspennstate/fourier-entropy-influence-one-coordinate}, 2026.
\bibitem{Shm3} Y.~Shmalo, Good coordinates in the Fourier entropy-influence conjecture,
\url{https://github.com/yspennstate/fei-good-coordinates}, 2026.
\bibitem{Shm6} Y.~Shmalo, Erasure entropies of sparse linear DNFs,
\url{https://github.com/yspennstate/fei-sparse-linear-dnf}, 2026.
\bibitem{WWW} A.~Wan, J.~Wright, C.~Wu, Decision trees, protocols and the entropy-influence conjecture, ITCS 2014,
arXiv:1312.3003.
\bibitem{MathBrainRepo} Y.~Shmalo, \emph{Math Brain: public snapshot, code and tests}, 2026,
\url{https://github.com/yspennstate/math-brain-public}.
\bibitem{MathBrainPaper} Y.~Shmalo, \emph{Math Brain: shared, auditable memory for multi-agent mathematical research},
October 2026, \url{https://github.com/yspennstate/math-brain-public/tree/main/papers/math-brain-2026-10-24}.
\bibitem{BrainFramework} Y.~Shmalo, \emph{Brain Framework}, public software repository, 2026,
\url{https://github.com/yspennstate/brain-framework-public}.
\end{thebibliography}

\end{document}
